
\subsection{Curvature away from the saddle on the antisymmetric branch}
\label{subsec:curvature-away-from-saddle-on-antisymmetric}
This subsection bounds the spherical product pushforward second variation at
anti-symmetric branch points away from the antipodal target.  On a fixed
metric annulus, sufficiently strong regularization dominates the fitting curvature.

We first note that the antipodal Dirac target lies exactly at the active
odd-moment origin by \cref{prop:intersection-zero-output-fiber}:
\[
    \mathbf S(\mu_\star^\pm)=0,
    \qquad
    \mathbf A(\mu_\star^\pm)=0.
\]

The fitting term may have negative curvature, but on bounded anti-symmetric
active sets this negative part is uniformly controlled; no metric annulus is
needed for the following fitting estimate.

\begin{lemma}[Fitting-curvature lower bound on anti-symmetric active balls]
\label{lem:fitting-curvature-lower-bound-Mminus-smoothness}
\textbf{Assumptions used.}
\cref{assumption:on-sphere,assumption:odd-symmetric-teacher},
together with the finite active odd-block representation of the fitting term.

Fix \(R_{\mathrm{act}}<\infty\).  Consider anti-symmetric branch points
\[
    \mu=\mu_1\otimes\mu_2\in\submanifoldforsaddlepoint,
    \qquad
    \Theta(\mu)=(0,\mathbf A),
    \qquad
    \|\mathbf A\|_{\mathcal H_{\mathcal O}}\le R_{\mathrm{act}}.
\]
This active-radius condition can always be met by taking
\(R_{\mathrm{act}}=\sqrt{|\mathcal O|}\).  Indeed,
\(\|T_{i,n}\|_{\mathrm{HS}}\le
\int\|w^{\otimes n}\|_{\mathrm{HS}}\,\dd\mu_i(w)=1\), so
\[
 \|A_n\|_{\mathrm{HS}}
 \le\frac{\|T_{1,n}\|_{\mathrm{HS}}+\|T_{2,n}\|_{\mathrm{HS}}}{2}\le1,
 \qquad
 \|\mathbf A\|_{\mathcal H_{\mathcal O}}^2
 =\sum_{n\in\mathcal O}\|A_n\|_{\mathrm{HS}}^2\le|\mathcal O|.
\]
Let \(v=(v_1,v_2)\) be arbitrary smooth spherical tangent perturbations,
\[
    v_i(w)\in T_w\Sphere^{d_w-1},
    \qquad
    w\cdot v_i(w)=0.
\]
The base point is on the anti-symmetric branch, but \(v\) is an arbitrary
ambient product tangent field and need not be tangent to that branch.
Let
\[
    R_{i,h}(w)
    \coloneq
    \frac{w+h v_i(w)}{\|w+h v_i(w)\|},
    \qquad
    \mu_{i,h}\coloneq (R_{i,h})_\#\mu_i,
    \qquad
    \mu_h\coloneq \mu_{1,h}\otimes\mu_{2,h}.
\]
Then there exists a constant
\[
    C_{\mathrm{fit}}^{-}(R_{\mathrm{act}})<\infty
\]
depending only on \(R_{\mathrm{act}}\), the finite active block \(\mathcal O\), the Hermite coefficients, and the teacher block, such that
\begin{equation}
\label{eq:fitting-curvature-lower-bound-Mminus-derived}
    \frac{\dd^2}{\dd h^2}\Big|_{h=0}
    F_{\mathrm{fit}}(\mu_h)
    \ge
    -C_{\mathrm{fit}}^{-}(R_{\mathrm{act}})
    \|v\|_{L^2(\mu)}^2,
\end{equation}
where
\[
    \|v\|_{L^2(\mu)}^2
    \coloneq
    \|v_1\|_{L^2(\mu_1)}^2
    +
    \|v_2\|_{L^2(\mu_2)}^2.
\]
Consequently, the abstract lower bound
\[
    \frac{\dd^2}{\dd h^2}\Big|_{h=0}
    F_{\mathrm{fit}}\bigl((R_{1,h})_\#\mu_1\otimes(R_{2,h})_\#\mu_2\bigr)
    \ge
    -C_{\mathrm{fit}}\|v\|_{L^2(\mu;\mathbb R^{d_w})}^2
\]
holds on the anti-symmetric active ball
\(\|\mathbf A\|_{\mathcal H_{\mathcal O}}\le R_{\mathrm{act}}\), with
\[
    C_{\mathrm{fit}}
    =
    C_{\mathrm{fit}}^{-}(R_{\mathrm{act}}).
\]
\end{lemma}

\begin{proof}
\noindent\textbf{Step 1: write the finite-moment second-derivative formula.}
The proof reduces the second variation to a finite-dimensional chain rule in
\(\Theta=(\mathbf S,\mathbf A)\).  Recall the finite active moment map from
\cref{eq:mf-finite-active-moment-map} and the ambient fitting polynomial from
\cref{eq:ambient-active-fit-functional}.
Since the fitting term is represented by the finite active odd block, we have
\[
    F_{\mathrm{fit}}(\mu_h)
    =
    \widetilde{\mathcal J}_{\mathrm{fit}}\bigl(\Theta(\mu_h)\bigr),
\]
where
\[
    \Theta(\mu_h)
    =
    \bigl(\mathbf S(h),\mathbf A(h)\bigr)
    \in
    \mathcal H_{\mathcal O}\times\mathcal H_{\mathcal O}.
\]
At \(h=0\), since \(\mu\in\submanifoldforsaddlepoint\),
\[
    \Theta(\mu)=(0,\mathbf A).
\]
For \(h\ne0\), an ambient perturbation may leave the branch, so both
components \(\mathbf S(h)\) and \(\mathbf A(h)\) are retained.
Set
\[
    \dot\Theta[v]
    \coloneq
    \frac{\dd}{\dd h}\Big|_{h=0}\Theta(\mu_h),
    \qquad
    \ddot\Theta[v]
    \coloneq
    \frac{\dd^2}{\dd h^2}\Big|_{h=0}\Theta(\mu_h).
\]
The ordinary chain rule gives
\begin{align}
\label{eq:fit-second-variation-chain-Mminus}
    \frac{\dd^2}{\dd h^2}\Big|_{h=0}
    F_{\mathrm{fit}}(\mu_h)
    =
    &\,
    D^2\widetilde{\mathcal J}_{\mathrm{fit}}(0,\mathbf A)
    [\dot\Theta[v],\dot\Theta[v]]
    \\
    &+
    \innerProd{
        \nabla\widetilde{\mathcal J}_{\mathrm{fit}}(0,\mathbf A)
    }{
        \ddot\Theta[v]
    }_{\mathcal H_{\mathcal O}\times\mathcal H_{\mathcal O}}.
\end{align}
The second term is required because the moment-coordinate path
\(h\mapsto\Theta(\mu_h)\) generally has the nonzero acceleration
\(\ddot\Theta[v]\).


\par\noindent\textbf{Step 2: bound the first and second moment derivatives.}
The fields use the norms of \(L^2(\mu_i;\mathbb R^{d_w})\); moment tensors
use their Hilbert--Schmidt norms; and \(\dot\Theta,\ddot\Theta\), and the
fitting gradient use the product norm of
\(\mathcal H_{\mathcal O}\times\mathcal H_{\mathcal O}\).
For each \(n\in\mathcal O\),
\[
    T_{i,n}(h)
    =
    \int R_{i,h}(w)^{\otimes n}\,\dd\mu_i(w).
\]
Tangency \(w^\top v_i(w)=0\) gives
\(\|w+hv_i(w)\|^2=1+h^2\|v_i(w)\|^2\). Expanding its inverse square
root in the normalized retraction yields
\begin{equation}
\label{eq:mf-retraction-second-jet}
\begin{aligned}
 R_{i,h}(w)&=\frac{w+hv_i(w)}{\sqrt{1+h^2\|v_i(w)\|^2}}
 =w+hv_i(w)-\tfrac12h^2\|v_i(w)\|^2w+O(h^3),\\
 \dot R_{i,0}(w)&=v_i(w),\qquad
 \ddot R_{i,0}(w)=-\|v_i(w)\|^2w,\qquad
 (I-ww^\top)\ddot R_{i,0}(w)=0.
\end{aligned}
\end{equation}
These derivatives agree with those of \(h\mapsto\EXP_w(hv_i(w))\)
from \cref{eq:mf-spherical-EXP}; the two curves need not agree beyond
second order. Applying the product rule to \(R_{i,h}(w)^{\otimes n}\),
we obtain
\begin{align}
    \dot T_{i,n}[v_i]
    &=
    n\int
    \sym{w^{\otimes(n-1)}\otimes v_i(w)}
    \,\dd\mu_i(w),
    \\
    \ddot T_{i,n}[v_i]
    &=
    n(n-1)\int
    \sym{w^{\otimes(n-2)}
        \otimes v_i(w)\otimes v_i(w)}
    \,\dd\mu_i(w)
    \\
    &\qquad
    -
    n\int
    \|v_i(w)\|^2 w^{\otimes n}
    \,\dd\mu_i(w).
\end{align}
The term containing \(n(n-1)\) is omitted when \(n=1\).
Because \(\|w\|=1\), Cauchy--Schwarz yields
\begin{equation}
\label{eq:dot-T-bound-fit-smoothness}
    \|\dot T_{i,n}[v_i]\|_{\HS}
    \le
    n\|v_i\|_{L^2(\mu_i)},
\end{equation}
and
\begin{equation}
\label{eq:ddot-T-bound-fit-smoothness}
    \|\ddot T_{i,n}[v_i]\|_{\HS}
    \le
    n^2\|v_i\|_{L^2(\mu_i)}^2.
\end{equation}
Recall the active coordinates \(\Theta=(\mathbf S,\mathbf A)\), where
\(S_n=(T_{1,n}+T_{2,n})/2\) and \(A_n=(T_{1,n}-T_{2,n})/2\).
The parallelogram identity gives
\[
 \|\dot\Theta\|^2=\frac12\sum_{n\in\mathcal O}
       (\|\dot T_{1,n}\|_{\HS}^2+\|\dot T_{2,n}\|_{\HS}^2),
 \qquad
 \|\ddot\Theta\|^2=\frac12\sum_{n\in\mathcal O}
       (\|\ddot T_{1,n}\|_{\HS}^2+\|\ddot T_{2,n}\|_{\HS}^2).
\]
Here both norms on \(\Theta\) are in
\(\mathcal H_{\mathcal O}\times\mathcal H_{\mathcal O}\).
Using \cref{eq:dot-T-bound-fit-smoothness,eq:ddot-T-bound-fit-smoothness}
and \(a^2+b^2\le(a+b)^2\) for nonnegative \(a,b\), we obtain
\begin{equation}
\label{eq:dot-Theta-bound-fit-smoothness}
 \|\dot\Theta[v]\|\le
 \left(\sum_{n\in\mathcal O}n^2\right)^{1/2}\|v\|_{L^2(\mu)},
\end{equation}
\begin{equation}
\label{eq:ddot-Theta-bound-fit-smoothness}
 \|\ddot\Theta[v]\|\le
 \left(\sum_{n\in\mathcal O}n^4\right)^{1/2}\|v\|_{L^2(\mu)}^2.
\end{equation}
For example, the square of the second left-hand side is at most
\(\tfrac12\sum_n n^4(\|v_1\|_{L^2(\mu_1)}^4+
\|v_2\|_{L^2(\mu_2)}^4)\), which is at most
\(\sum_n n^4\|v\|_{L^2(\mu)}^4\).


\par\noindent\textbf{Step 3: choose the curvature bound and apply the chain rule.}
The set \(\mathcal O\) is finite and each tensor block is finite-dimensional.
Hence
\(\{(0,\mathbf A):\|\mathbf A\|_{\mathcal H_{\mathcal O}}
\le R_{\mathrm{act}}\}\) is compact.  The fitting functional
\(\widetilde{\mathcal J}_{\mathrm{fit}}\) is a polynomial, so its first two
derivatives are continuous and have finite suprema on this set.
Define the fitting curvature constant directly by
\begin{equation}
\label{eq:def-Cfit-Mminus-derived}
\begin{aligned}
 C_{\mathrm{fit}}^-(R_{\mathrm{act}}):={}&
 \left(\sum_{n\in\mathcal O}n^2\right)
 \sup_{\|\mathbf A\|\le R_{\mathrm{act}}}
       \|D^2\widetilde{\mathcal J}_{\mathrm{fit}}(0,\mathbf A)\|_{\mathrm{op}}\\
 &+\left(\sum_{n\in\mathcal O}n^4\right)^{1/2}
 \sup_{\|\mathbf A\|\le R_{\mathrm{act}}}
       \|\nabla\widetilde{\mathcal J}_{\mathrm{fit}}(0,\mathbf A)\|.
\end{aligned}
\end{equation}
The gradient norm is the product \(\mathcal H_{\mathcal O}\) norm.
The Hessian norm in \cref{eq:def-Cfit-Mminus-derived} is the operator norm of
\[
 D^2\widetilde{\mathcal J}_{\mathrm{fit}}(0,\mathbf A):
 \mathcal H_{\mathcal O}\times\mathcal H_{\mathcal O}
 \longrightarrow
 \mathcal H_{\mathcal O}\times\mathcal H_{\mathcal O}.
\]
The chain rule \cref{eq:fit-second-variation-chain-Mminus} and
Cauchy--Schwarz now give
\[
\begin{aligned}
 \left.\frac{\dd^2}{\dd h^2}F_{\mathrm{fit}}(\mu_h)\right|_0
 &\ge-\|D^2\widetilde{\mathcal J}_{\mathrm{fit}}(0,\mathbf A)\|_{\mathrm{op}}
                   \|\dot\Theta\|^2
       -\|\nabla\widetilde{\mathcal J}_{\mathrm{fit}}(0,\mathbf A)\|
                   \|\ddot\Theta\|\\
 &\ge-C_{\mathrm{fit}}^-(R_{\mathrm{act}})\|v\|_{L^2(\mu)}^2.
\end{aligned}
\]
This proves \cref{eq:fitting-curvature-lower-bound-Mminus-derived} with
the constant defined in \cref{eq:def-Cfit-Mminus-derived}.
\end{proof}

The metric regularizer relies on the positive Hessian of squared spherical distance inside each fixed cap.
\begin{lemma}[Spherical Hessian of the fixed squared-distance cost]
\label{lem:spherical-squared-distance-Hessian-on-caps}
\textbf{Assumptions used.} \cref{assumption:on-sphere}.

Let
\[
    c_y(w)=\frac12d_\Sphere(w,y)^2,
    \qquad
    y\in\Sphere^{d_w-1}.
\]
On the geodesic ball \(d_\Sphere(w,y)\le r_{\mathrm{cap}}<\pi/2\), the spherical Hessian satisfies
\begin{equation}
\label{eq:spherical-distance-Hessian-lower-bound}
    \nabla_\Sphere^2c_y(w)[v,v]
    \ge
    \kappa_{\mathrm{cap}}\|v\|^2,
    \qquad
    \kappa_{\mathrm{cap}}
    \coloneq
    r_{\mathrm{cap}}\cot r_{\mathrm{cap}}>0.
\end{equation}
Consequently, for \(c_\star(w)=\frac12d_\Sphere(w,\pi_\star(w))^2\),
\begin{equation}
\label{eq:cstar-Hessian-lower-bound}
    \nabla_\Sphere^2c_\star(w)[v,v]
    \ge
    \kappa_{\mathrm{cap}}
    \|v\|^2
\end{equation}
for every \(w\in K_+(r_{\mathrm{cap}})\cup K_-(r_{\mathrm{cap}})\) and every
\(v\in T_w\Sphere^{d_w-1}\).
\end{lemma}

\begin{proof}
For \(0<r=d_\Sphere(w,y)<\pi/2\), define
\(e=(y-\cos r\,w)/\sin r\).  Since \(w^\top y=\cos r\),
\[
 w^\top e=\frac{w^\top y-\cos r}{\sin r}=0,
 \qquad
 \|e\|_{\mathbb R^{d_w}}^2
 =\frac{1-2\cos r\,y^\top w+\cos^2r}{\sin^2r}=1.
\]
Thus \(e\in T_w\Sphere^{d_w-1}\) is the unit tangent pointing from \(w\)
toward \(y\).
Put \(x(t)=y^\top\EXP_w(tv)\) and \(f(x)=\tfrac12(\arccos x)^2\).
By \cref{eq:mf-spherical-EXP},
\[
 x(0)=\cos r,\quad x'(0)=y^\top v=\sin r\,\langle e,v\rangle,
 \quad x''(0)=-\|v\|^2\cos r,
\]
and direct differentiation gives
\[
 f'(x)=-\frac{\arccos x}{\sqrt{1-x^2}},\qquad
 f''(x)=\frac1{1-x^2}
        -\frac{x\arccos x}{(1-x^2)^{3/2}}.
\]
Substituting \(x=\cos r\) yields
\[
 f'(\cos r)=-\frac r{\sin r},\qquad
 f''(\cos r)=\frac{\sin r-r\cos r}{\sin^3r}.
\]
For sufficiently small \(|t|\), the EXP path does not reach the antipode of
\(y\), so the principal \(\arccos\) gives the spherical distance.  Hence
\[
 c_y(\EXP_w(tv))
 =\frac12d_\Sphere(\EXP_w(tv),y)^2
 =\frac12[\arccos(y^\top\EXP_w(tv))]^2
 =f(x(t)).
\]
The one-variable chain rule now gives
\begin{equation}
\label{eq:mf-spherical-cost-Hessian-formula}
\begin{aligned}
 \nabla_\Sphere^2c_y(w)[v,v]
 &=f''(\cos r)(x'(0))^2+f'(\cos r)x''(0)\\
 &=r\cot r\,\|v\|^2+(1-r\cot r)\langle v,e\rangle^2.
\end{aligned}
\end{equation}
To display the eigendirections, set
\(v_\parallel=\langle v,e\rangle e\) and
\(v_\perp=v-v_\parallel\).  Since
\(\|v\|_{\mathbb R^{d_w}}^2=\|v_\parallel\|_{\mathbb R^{d_w}}^2+
\|v_\perp\|_{\mathbb R^{d_w}}^2\),
\[
 \nabla_\Sphere^2c_y(w)[v,v]
 =\|v_\parallel\|_{\mathbb R^{d_w}}^2
  +r\cot r\,\|v_\perp\|_{\mathbb R^{d_w}}^2.
\]
The radial eigenvalue is therefore \(1\), and the orthogonal eigenvalue is
\(r\cot r\). Moreover,
\[
 (r\cot r)'=\frac{\sin r\cos r-r}{\sin^2r}\le0,\qquad
 r-\sin r\cos r=\int_0^r2\sin^2s\,\dd s\ge0,
 \qquad \lim_{r\to0}r\cot r=1.
\]
The integral identity shows that \((r\cot r)'\le0\).  Therefore
\(1\ge r\cot r\ge r_{\mathrm{cap}}\cot r_{\mathrm{cap}}>0\) for
\(0<r\le r_{\mathrm{cap}}\), so both eigenvalues are at least
\(\kappa_{\mathrm{cap}}\).
At \(w=y\), the identity
\(c_y(\EXP_y(tv))=t^2\|v\|_{\mathbb R^{d_w}}^2/2\) for small \(t\)
gives Hessian \(\|v\|_{\mathbb R^{d_w}}^2\) without introducing \(e\).
On each cap, \(\pi_\star(w)\) is the fixed point \(a_\star\) or
\(-a_\star\), so the calculation does not differentiate across a cap
boundary.  Substituting
\(y=a_\star\) or \(y=-a_\star\) gives the claim for the fixed-assignment
cost \(c_\star(w)=d_\Sphere(w,\pi_\star(w))^2/2\)
from \cref{eq:mf-fixed-assignment-cost,eq:mf-cap-class}.
\end{proof}

Combining fixed assignment with the spherical distance Hessian gives a positive spherical product pushforward second variation for the metric regularizer.
\begin{lemma}[Metric-power regularizer Hessian on fixed-assignment caps]
\label{lem:W2-power-regularizer-Hessian-on-fixed-caps}
\textbf{Assumptions used.} \cref{assumption:on-sphere,assumption:odd-symmetric-teacher}, together with \cref{lemma-fixed-assignment,lem:spherical-squared-distance-Hessian-on-caps}.

Let \(\mu_i\in\mathcal P_\pm(r_{\mathrm{cap}})\), and let \(v=(v_1,v_2)\) be smooth tangent perturbations whose spherical retractions preserve the fixed-assignment caps for sufficiently small \(|h|\).  Then
\begin{align}
\label{eq:W2-power-Hessian-lower-bound-fixed-caps}
    D^2_{W_2}F_{\mathrm{reg}}(\mu_1,
    \mu_2)[v,v]
    \ge
    \lambdaReg\,
    \kappa_{\mathrm{cap}}
    Z_\pm(\mu_1,
    \mu_2)^{\delta_W/2}
    \|v\|_{L^2(\mu)}^2.
\end{align}
In particular, if \(Z_\pm(\mu_1,
\mu_2)\ge\zeta_{\min}^2\), then
\begin{equation}
\label{eq:W2-power-Hessian-lower-bound-annulus}
    D^2_{W_2}F_{\mathrm{reg}}(
    \mu_1,
    \mu_2)[v,v]
    \ge
    \lambdaReg\,
    \kappa_{\mathrm{cap}}
    \zeta_{\min}^{\delta_W}
    \|v\|_{L^2(\mu)}^2.
\end{equation}
At \(Z_\pm=0\), the regularizer has zero quadratic variation because \(p_W=2+\delta_W>2\).
\end{lemma}

\begin{proof}
\noindent\textbf{Step 1: bound the positive-distance second derivative.}
By \cref{eq:mf-retraction-second-jet}, the normalized retraction used here
and the spherical EXP path have the same first and second derivatives at
\(h=0\).  Since the regularizer is smooth when \(Z_\pm>0\), its second
derivative along the retraction is therefore the one computed for the EXP
path in \cref{eq:second-variation-W2-dirac-metric-regularizer}.  In that
formula, the second term is nonnegative because
\(p_W-2=\delta_W>0\).  Substituting the pointwise Hessian bound
\cref{eq:cstar-Hessian-lower-bound} into the first term gives
\begin{align*}
 D^2_{W_2}F_{\mathrm{reg}}(\mu_1,\mu_2)[v,v]
 &\ge
 \lambdaReg\kappa_{\mathrm{cap}}Z_\pm^{p_W/2-1}
 \sum_{i=1}^2\int
 \|v_i(w)\|_{\mathbb R^{d_w}}^2\,\dd\mu_i(w)\\
 &=
 \lambdaReg\kappa_{\mathrm{cap}}Z_\pm^{\delta_W/2}
 \|v\|_{L^2(\mu;\mathbb R^{d_w})}^2,
\end{align*}
which is \cref{eq:W2-power-Hessian-lower-bound-fixed-caps} for
\(Z_\pm>0\).  If \(Z_\pm\ge\zeta_{\min}^2\), then \(\delta_W>0\) implies
\[
 Z_\pm^{\delta_W/2}
 \ge (\zeta_{\min}^2)^{\delta_W/2}
 =\zeta_{\min}^{\delta_W},
\]
and \cref{eq:W2-power-Hessian-lower-bound-annulus} follows.

\par\noindent\textbf{Step 2: compute the ordinary second derivative at the target.}
For \(Z_\pm=0\), both marginals equal
\(\nu_\star^\pm=(\delta_{a_\star}+\delta_{-a_\star})/2\).
At a pole \(y=\varepsilon a_\star\), tangency gives
\(y^\top R_{i,h}(y)=(1+h^2\|v_i(y)\|^2)^{-1/2}\), so its distance
from the pole is \(\arctan(|h|\|v_i(y)\|)\). The fixed-assignment
formula \cref{eq:fixed-assignment-W2-square-formula} consequently gives
\begin{equation}
\label{eq:mf-zero-distance-retraction}
 Z_\pm(\mu_h)=\frac12\sum_{i=1}^2\sum_{\varepsilon\in\{-1,1\}}
 \arctan^2(h\|v_i(\varepsilon a_\star)\|)
 =h^2\|v\|_{L^2(\mu)}^2+O(h^4).
\end{equation}
If \(\|v\|_{L^2(\mu)}=0\), all four atoms stay fixed. Otherwise
\(a(h):=Z_\pm(\mu_h)/h^2\) extends to a smooth positive function at
zero with \(a(0)=\|v\|_{L^2(\mu)}^2\). Set
\(b(h):=(\lambdaReg/p_W)a(h)^{p_W/2}\), so the regularizer is
\(|h|^{p_W}b(h)\). For \(h\ne0\), differentiating twice gives
\[
\begin{aligned}
 F_{\mathrm{reg}}'(h)
 &=p_W\operatorname{sgn}(h)|h|^{p_W-1}b(h)+|h|^{p_W}b'(h),\\
 F_{\mathrm{reg}}''(h)
 &=p_W(p_W-1)|h|^{p_W-2}b(h)
 +2p_W\operatorname{sgn}(h)|h|^{p_W-1}b'(h)+|h|^{p_W}b''(h).
\end{aligned}
\]
The difference quotients at zero and these formulas show that the
regularizer is \(C^2\), with both derivatives zero there, because
\(p_W>2\). This proves the bound at the target as well.
\end{proof}

The next proposition is an abstract combination statement: at a fixed state,
it adds any available fitting lower bound to the metric lower bound.  On an
anti-symmetric active ball, the fitting input has already been verified in
\cref{lem:fitting-curvature-lower-bound-Mminus-smoothness}; the subsequent
corollary makes that substitution.

\begin{proposition}[Ambient second-variation bound on fixed metric annuli]
\label{prop:sqchi2-away-from-origin-Wasserstein-convexity}
\textbf{Assumptions used.} \cref{assumption:on-sphere,assumption:odd-symmetric-teacher}, together with \cref{lem:W2-power-regularizer-Hessian-on-fixed-caps}.

Fix \(0<\zeta_{\min}<\infty\), and consider a state in a fixed-assignment
cap region satisfying
\[
    \mu_i\in\mathcal P_\pm(r_{\mathrm{cap}}),
    \qquad
    Z_\pm(\mu_1,
    \mu_2)
    \ge
    \zeta_{\min}^2.
\]
Suppose that, on the region under consideration, the fitting lower bound
\begin{equation}
\label{eq:fit-curvature-lower-bound-abstract}
    \frac{\dd^2}{\dd h^2}\Big|_{h=0}
    F_{\mathrm{fit}}\bigl((R_{1,h})_\#\mu_1\otimes(R_{2,h})_\#\mu_2\bigr)
    \ge
    -C_{\mathrm{fit}}\|v\|_{L^2(\mu)}^2
\end{equation}
is available for every smooth tangent perturbation \(v=(v_1,v_2)\) whose
retractions stay in the fixed-assignment caps for small \(|h|\).  Then the
full objective satisfies the following pointwise bound at that state:
\begin{equation}
\label{eq:sqchi2-away-from-origin-W2-convexity-bound}
    \innerProd{v}{\mathcal H_{\Sphere,\mu}v}_{L^2(\mu;\mathbb R^{d_w})}
    \ge
    \left(
        \lambdaReg\kappa_{\mathrm{cap}}\zeta_{\min}^{\delta_W}
        -C_{\mathrm{fit}}
    \right)
    \|v\|_{L^2(\mu;\mathbb R^{d_w})}^2.
\end{equation}
In particular, if
\[
    \lambdaReg>
    \frac{C_{\mathrm{fit}}}{\kappa_{\mathrm{cap}}\zeta_{\min}^{\delta_W}},
\]
then the spherical product pushforward second variation is strictly positive in
every nonzero admissible tangent direction at each point of the region.
\end{proposition}

\begin{proof}
By \cref{lem:mf-full-spherical-Hessian}, the spherical product pushforward
second variation is
\[
 D_{W_2}^2F(\mu)[v,v]
 =\innerProd{v}{\mathcal H_{\Sphere,\mu}v}_{L^2(\mu;\mathbb R^{d_w})}.
\]
The two inputs give
\begin{align*}
 D_{W_2}^2F_{\mathrm{fit}}(\mu)[v,v]
 &\ge -C_{\mathrm{fit}}\|v\|_{L^2(\mu;\mathbb R^{d_w})}^2,
 &&\text{by \cref{eq:fit-curvature-lower-bound-abstract},}\\
 D_{W_2}^2F_{\mathrm{reg}}(\mu)[v,v]
 &\ge \lambdaReg\kappa_{\mathrm{cap}}\zeta_{\min}^{\delta_W}
        \|v\|_{L^2(\mu;\mathbb R^{d_w})}^2,
 &&\text{by \cref{eq:W2-power-Hessian-lower-bound-annulus}.}
\end{align*}
Adding these inequalities proves
\cref{eq:sqchi2-away-from-origin-W2-convexity-bound}.
\end{proof}

\begin{remark}[Pointwise scope]
The annulus need not be geodesically convex.  A strong-convexity inequality
between two distinct states follows only along a comparison geodesic that
remains in a region where
\cref{eq:sqchi2-away-from-origin-W2-convexity-bound} holds at every point.
\end{remark}

Specializing this bound to the anti-symmetric branch gives positive ambient
curvature at the prescribed fixed-assignment annular states.

\begin{corollary}[Ambient curvature at anti-symmetric annular states]
\label{cor:branchwise-positivity-Wasserstein-Mminus}
\textbf{Assumptions used.}
\cref{assumption:on-sphere,assumption:odd-symmetric-teacher}, together with
\cref{prop:sqchi2-away-from-origin-Wasserstein-convexity}.

Fix \(R_{\mathrm{act}}<\infty\).  Let \(\mu_-\in\submanifoldforsaddlepoint\)
satisfy \(\|\mathbf A(\mu_-)\|_{\mathcal H_{\mathcal O}}\le R_{\mathrm{act}}\) and
\[
    \mu_{-,i}\in\mathcal P_\pm(r_{\mathrm{cap}}),
    \qquad
    Z_\pm(\mu_{-,1},\mu_{-,2})
    \ge
    \zeta_{\min}^2.
\]
Then \cref{eq:sqchi2-away-from-origin-W2-convexity-bound} holds with
\(C_{\mathrm{fit}}=C_{\mathrm{fit}}^{-}(R_{\mathrm{act}})\) from
\cref{eq:fitting-curvature-lower-bound-Mminus-derived,eq:def-Cfit-Mminus-derived}.
In particular, every nonzero
admissible product tangent field has positive spherical product pushforward
second variation when
\(\lambdaReg\kappa_{\mathrm{cap}}\zeta_{\min}^{\delta_W}>
C_{\mathrm{fit}}^{-}(R_{\mathrm{act}})\).  The field need not be tangent
to \(\submanifoldforsaddlepoint\).
\end{corollary}
